\section{Einleitung}\label{einleitung}

Die Verbreitung von Internet of Things (IoT)-Geräten in privaten, industriellen und kritischen Infrastrukturen vergrößert die Angriffsfläche moderner Netzwerke kontinuierlich \cite{hamedani2025iotsurvey}. Hamedani et al. \cite{hamedani2025iotsurvey} zeigen in einem umfassenden Survey über ML-basierte Sicherheitslösungen für IoT, dass die Forschung in diesem Feld zwischen 2020 und 2024 stark gewachsen ist und gleichzeitig methodische Lücken bestehen, insbesondere in der Bewertung und der Interpretierbarkeit der eingesetzten Modelle. Kikissagbe und Adda \cite{kikissagbe2024review} kategorisieren in ihrem systematischen Review supervised, unsupervised und hybride ML-Ansätze für IoT-Intrusion Detection Systeme (IDS) und ordnen baumbasierte Ensemble-Methoden als etablierten Forschungsstand ein.

Eine zentrale empirische Studie auf diesem Forschungsfeld ist die Arbeit von Raturi et al. \cite{raturi2026exploratory}, die mehrere ML-Algorithmen auf dem CICIoT2023-Datensatz \cite{neto2023ciciot2023} evaluiert. Gradient Boosting erzielt in der Mehrklassenklassifikation die höchste Balanced Accuracy unter den getesteten Modellen \cite{raturi2026exploratory}. Die Autoren benennen jedoch drei methodische Lücken, die ihre Studie offen lässt: die fehlende Erklärbarkeit der Modellentscheidungen, die fehlende Analyse der Fehlerfortpflanzung im zweistufigen Klassifikationsansatz und das Fehlen von Latenzmessungen \cite{raturi2026exploratory}. Die vorliegende Arbeit adressiert die erste dieser Lücken.

Die Wahl von SHapley Additive Explanations (SHAP) als Erklärbarkeitsmethode ist wissenschaftlich begründet. Ogunseyi et al. \cite{ogunseyi2026xai_review} analysieren in einem PRISMA-basierten systematischen Review 129 Studien aus den Jahren 2018 bis 2025 zur explainable Artificial Intelligence (XAI) in IDS für IoT-Netzwerke und zeigen, dass SHAP und Local Interpretable Model-agnostic Explanations (LIME) in der überwiegenden Mehrheit der untersuchten Studien eingesetzt werden, ohne die Erkennungsgenauigkeit der zugrundeliegenden Modelle zu beeinträchtigen. Mohale und Obagbuwa \cite{mohale2025xai} demonstrieren am Beispiel mehrerer ML-Klassifikatoren einschließlich XGBoost, wie globale und lokale SHAP-Analysen in IDS-Experimenten umgesetzt werden können. Hermosilla, Berríos und Allende-Cid \cite{hermosilla2025xai_forensic} ergänzen diesen methodischen Rahmen um quantitative Bewertungsmetriken für die Qualität der Erklärungen, darunter Fidelität, Konsistenz und Jaccard-Ähnlichkeit, evaluiert auf XGBoost und TabNet im Kontext der Intrusion Detection.

Die vorliegende Arbeit verbindet diese methodischen Bausteine. Sie entwickelt einen Gradient-Boosting-Klassifikator auf dem CICIoT2023-Datensatz, evaluiert ihn mit der Balanced Accuracy als primärer Metrik und integriert SHAP als post-hoc-Erklärungskomponente. Eine Ablation Study untersucht zusätzlich den Einfluss der Principal Component Analysis (PCA) auf Modellleistung und Erklärungen.

Die Arbeit ist wie folgt strukturiert. Kapitel~2 ordnet die zehn verwendeten Quellen thematisch in vier Cluster ein und arbeitet die Forschungslücke heraus. Kapitel~3 beschreibt den CICIoT2023-Datensatz, dessen Klassenungleichgewicht und die geplante Vorverarbeitung einschließlich der zwei parallelen Preprocessing-Pipelines. Kapitel~4 dokumentiert die methodische Vorgehensweise in fünf Teilaspekten: die zweistufige hierarchische Klassifikationsarchitektur nach Raturi et al. \cite{raturi2026exploratory}, die Wahl von Gradient Boosting als Hauptmodell, die Balanced Accuracy als primäre Evaluationsmetrik begründet durch Hosseini et al. \cite{hosseini2025imbalance}, die Ablation Study zum PCA-Einfluss auf Modellleistung und SHAP-Erklärungen, sowie die SHAP-Integration mit den Qualitätsmetriken aus Hermosilla et al. \cite{hermosilla2025xai_forensic}. Kapitel~5 dokumentiert die Implementierung. Kapitel~6 präsentiert die Evaluationsergebnisse beider Pipelines. Kapitel~7 diskutiert die Ergebnisse im Vergleich zur bestehenden Literatur, insbesondere gegenüber Raturi et al. \cite{raturi2026exploratory} und Alharby \cite{alharby2025ciciot}. Kapitel~8 schließt mit Fazit und Ausblick.
